\documentclass[12pt,a4paper]{article}
\usepackage[utf8]{inputenc}
\usepackage[T2A]{fontenc}
\usepackage[russian]{babel}
\usepackage{amsmath,amssymb}
\usepackage[left=25mm,right=25mm,top=20mm,bottom=20mm]{geometry}
\usepackage{enumitem}

\title{Простые модели охлаждения нефти при добыче и разгазировании}
\author{}
\date{}

\begin{document}
\maketitle

\section{Введение}

При добыче нефти наблюдается снижение температуры продукции, обусловленное
двумя основными механизмами:
\begin{itemize}[nosep]
  \item дросселирование потока при снижении давления (эффект Джоуля--Томсона);
  \item выделение растворённого газа из нефти (разгазирование).
\end{itemize}

Ниже представлены три уровня моделирования этих эффектов --- от простейшей
оценки до энтальпийного баланса двухфазной смеси. Все модели существенно
проще полной модели <<black oil>>, но сохраняют ключевую физику процесса.

%=============================================================================
\section{Модель~1. Эффект Джоуля--Томсона (дросселирование)}
%=============================================================================

Простейшая оценка: при адиабатическом снижении давления температура
однофазного флюида изменяется по формуле
\begin{equation}\label{eq:JT}
  \Delta T = \mu_{\mathrm{JT}}\,\Delta P,
\end{equation}
где $\mu_{\mathrm{JT}}$~--- коэффициент Джоуля--Томсона
(${}^\circ\mathrm{C}/\text{атм}$), $\Delta P$~--- перепад давления.

Типичные значения:
\begin{center}
\begin{tabular}{lc}
  \hline
  Флюид & $\mu_{\mathrm{JT}}$, ${}^\circ\mathrm{C}/\text{атм}$ \\
  \hline
  Нефть (дегазированная) & $0{,}03 \ldots 0{,}05$ \\
  Попутный газ           & $0{,}3  \ldots 0{,}5$  \\
  Вода                   & $\sim 0{,}02$           \\
  \hline
\end{tabular}
\end{center}

\paragraph{Область применимости.}
Формула~\eqref{eq:JT} корректна для однофазного потока и даёт оценку
охлаждения на штуцерах, в клапанах и при депрессии на пласт, когда
разгазирования не происходит ($P > P_{\text{нас}}$).

%=============================================================================
\section{Модель~2. Энергобаланс при разгазировании}
%=============================================================================

При снижении давления ниже давления насыщения $P_{\text{нас}}$ из нефти
выделяется растворённый газ. Процесс требует затрат энергии на:
\begin{enumerate}[nosep]
  \item разрыв межмолекулярных связей (<<скрытая теплота>> дегазации~$L$);
  \item работу расширения газовой фазы.
\end{enumerate}

Запишем баланс энергии для адиабатического процесса. Пусть при снижении
давления от $P_1$ до $P_2$ из единицы массы живой нефти выделяется масса газа
$\Delta m_g$. Тогда
\begin{equation}\label{eq:flash}
  \boxed{\;
    \Delta T = -\,\frac{L\,\Delta m_g
      + \dfrac{\Delta m_g}{\rho_g}\,(P_1 - P_2)}
    {m_o\,c_{p,o} + m_g\,c_{p,g}}
  \;}
\end{equation}
где:
\begin{itemize}[nosep]
  \item $L$~--- удельная теплота дегазации, Дж/кг
        (типично $150\ldots350$~кДж/кг, зависит от состава газа);
  \item $\Delta m_g$~--- масса выделившегося газа, кг;
  \item $\rho_g$~--- плотность газа при средних $P$, $T$;
  \item $m_o$, $m_g$~--- массы нефтяной и газовой фаз;
  \item $c_{p,o}$, $c_{p,g}$~--- удельные теплоёмкости фаз.
\end{itemize}

Первое слагаемое в числителе~--- затраты на дегазацию, второе~--- работа
расширения газа.

\paragraph{Связь с газовым фактором.}
Масса выделившегося газа определяется через газосодержание (газовый фактор):
\begin{equation}\label{eq:Rs}
  \Delta m_g = \rho_{g,\text{ст}}\,\bigl[R_s(P_1,T_1) - R_s(P_2,T_2)\bigr],
\end{equation}
где $R_s(P,T)$~--- удельное газосодержание (м$^3$/м$^3$ при стандартных условиях).
Для его оценки можно использовать корреляцию Стендинга:
\begin{equation}\label{eq:standing}
  R_s = \gamma_g\,\left[\left(\frac{P}{a}\right)
        \cdot 10^{\,b\,\gamma_{\mathrm{API}} - c\,T}\right]^d,
\end{equation}
где $\gamma_g$~--- относительная плотность газа, $\gamma_{\mathrm{API}}$~---
плотность нефти в градусах~API, $a$, $b$, $c$, $d$~--- эмпирические
константы ($a=18{,}2$, $b=0{,}0125$, $c=0{,}00091$, $d=1{,}2048$ в системе
единиц с давлением в psia и температурой в~${}^\circ\mathrm{F}$).

%=============================================================================
\section{Модель~3. Энтальпийный баланс двухфазного потока}
%=============================================================================

\subsection{Постановка задачи}

Рассматривается адиабатический стационарный поток нефтегазовой смеси
между двумя сечениями (<<1>> и <<2>>). Общая масса смеси $M$~(кг) сохраняется.
В сечении~1 флюид находится при давлении~$P_1$ и температуре~$T_1$;
в сечении~2 давление снижается до~$P_2 < P_1$.

Требуется найти температуру~$T_2$ в сечении~2.

\paragraph{Допущения.}
\begin{enumerate}[nosep]
  \item Процесс адиабатический (нет теплообмена с окружающей средой).
  \item Кинетической и потенциальной энергией потока пренебрегаем.
  \item Фазы находятся в термодинамическом равновесии (мгновенная сепарация).
  \item Свойства фаз ($c_p$, $\rho$) зависят от~$P$ и~$T$, но вычисляются
        по эмпирическим корреляциям (не из уравнения состояния).
\end{enumerate}

%-----------------------------------------------------------------------------
\subsection{Обозначения}
%-----------------------------------------------------------------------------

Во всех формулах ниже:
\begin{center}
\renewcommand{\arraystretch}{1.25}
\begin{tabular}{cl}
  \hline
  Символ & Смысл \\
  \hline
  $P$, $T$                  & давление (Па) и температура (К) \\
  индексы $1$, $2$          & сечение~1 (вход) и сечение~2 (выход) \\
  индекс $o$                & нефтяная (жидкая) фаза \\
  индекс $g$                & газовая фаза \\
  индекс $\text{ст}$        & стандартные условия: $P_{\text{ст}} = 0{,}1013$~МПа,
                               $T_{\text{ст}} = 20\,{}^\circ$С (РФ) или $15{,}6\,{}^\circ$С ($60\,{}^\circ$F, API) \\
  индекс $\mathrm{ref}$     & условия отсчёта ($P_{\mathrm{ref}}$, $T_{\mathrm{ref}}$) \\
  $M$                       & полная масса смеси, кг \\
  $m_o$, $m_g$              & массы жидкой и газовой фаз, кг; $m_o + m_g = M$ \\
  $x_g$                     & массовая доля газовой фазы: $x_g = m_g / M$ \\
  $R_s(P,T)$                & газосодержание, м$^3_{\text{ст}}$/м$^3_{\text{ст}}$ нефти \\
  $\rho_{o}(P,T)$           & плотность нефти при $(P,T)$, кг/м$^3$ \\
  $\rho_{g}(P,T)$           & плотность газа при $(P,T)$, кг/м$^3$ \\
  $\rho_{o,\text{ст}}$, $\rho_{g,\text{ст}}$
                             & плотности при стандартных условиях, кг/м$^3$ \\
  $c_{p,o}$, $c_{p,g}$      & удельные теплоёмкости при пост.\ давлении, Дж/(кг$\cdot$К) \\
  $L$                        & удельная теплота испарения (дегазации), Дж/кг газа \\
  $v_o = 1/\rho_o$, $v_g = 1/\rho_g$ & удельные объёмы, м$^3$/кг \\
  $\alpha_o$, $\alpha_g$    & коэффициенты теплового расширения, К$^{-1}$ \\
  $h_o$, $h_g$              & удельные энтальпии фаз, Дж/кг \\
  $h_{\mathrm{mix}}$        & удельная энтальпия смеси, Дж/кг \\
  \hline
\end{tabular}
\end{center}

%-----------------------------------------------------------------------------
\subsection{Закон сохранения энергии}
%-----------------------------------------------------------------------------

Из первого начала термодинамики для стационарного адиабатического потока
без изменения кинетической энергии следует сохранение удельной энтальпии
смеси:
\begin{equation}\label{eq:enthalpy}
  \boxed{\;h_{\mathrm{mix}}(P_1, T_1) = h_{\mathrm{mix}}(P_2, T_2).\;}
\end{equation}

%-----------------------------------------------------------------------------
\subsection{Энтальпия смеси}
%-----------------------------------------------------------------------------

Удельная энтальпия смеси~--- средневзвешенная по массе:
\begin{equation}\label{eq:Hmix}
  h_{\mathrm{mix}}(P,T) = \bigl[1 - x_g(P,T)\bigr]\,h_o(P,T)
                         + x_g(P,T)\,h_g(P,T).
\end{equation}

%-----------------------------------------------------------------------------
\subsection{Массовая доля газовой фазы $x_g$}
%-----------------------------------------------------------------------------

Газосодержание $R_s(P,T)$ определяет объём газа (при стандартных условиях),
растворённого в единице объёма нефти (при стандартных условиях). Если
при давлении насыщения $P_{\text{нас}}$ всё количество газа растворено,
а при давлении~$P < P_{\text{нас}}$ часть газа выделяется, то масса
свободного газа на единицу массы дегазированной нефти:
\begin{equation}\label{eq:mg_spec}
  \frac{m_g}{m_{o,\text{ст}}}
    = \frac{\rho_{g,\text{ст}}}{\rho_{o,\text{ст}}}\,
      \bigl[R_{s,\max} - R_s(P,T)\bigr],
\end{equation}
где $R_{s,\max} = R_s(P_{\text{нас}}, T)$~--- максимальное газосодержание,
$m_{o,\text{ст}}$~--- масса дегазированной нефти.

Полная масса смеси на единицу массы дегазированной нефти:
\begin{equation}\label{eq:M_spec}
  \frac{M}{m_{o,\text{ст}}}
    = 1 + \frac{\rho_{g,\text{ст}}}{\rho_{o,\text{ст}}}\,R_{s,\max}.
\end{equation}
Тогда массовая доля свободного газа:
\begin{equation}\label{eq:xg}
  \boxed{\;
    x_g(P,T) = \frac{m_g}{M}
    = \frac{\rho_{g,\text{ст}}\,\bigl[R_{s,\max} - R_s(P,T)\bigr]}
           {\rho_{o,\text{ст}} + \rho_{g,\text{ст}}\,R_{s,\max}}.
  \;}
\end{equation}

Заметим: при $P \geqslant P_{\text{нас}}$ имеем $R_s = R_{s,\max}$, и
$x_g = 0$ (весь газ растворён, однофазный поток).

%-----------------------------------------------------------------------------
\subsection{Удельные энтальпии фаз}
%-----------------------------------------------------------------------------

Энтальпия есть функция состояния. Для каждой фазы запишем
дифференциал:
\begin{equation}\label{eq:dh_general}
  dh = c_p\,dT + v\,(1 - T\,\alpha)\,dP,
\end{equation}
где $v = 1/\rho$~--- удельный объём, $\alpha = -\frac{1}{\rho}
\left(\frac{\partial\rho}{\partial T}\right)_P$~--- коэффициент теплового
расширения. Соотношение~\eqref{eq:dh_general} следует из определения
энтальпии и соотношения Максвелла $\left(\frac{\partial v}{\partial T}
\right)_P = v\,\alpha$.

Интегрируя~\eqref{eq:dh_general} от условий отсчёта
$(P_{\mathrm{ref}}, T_{\mathrm{ref}})$ до $(P,T)$:

\paragraph{Нефтяная фаза.}
\begin{equation}\label{eq:Ho}
  h_o(P,T) = c_{p,o}\,(T - T_{\mathrm{ref}})
            + v_o\,(1 - T\,\alpha_o)\,(P - P_{\mathrm{ref}}).
\end{equation}

\paragraph{Газовая фаза.}
Энтальпию газа отсчитываем от жидкой фазы, прибавляя теплоту испарения~$L$:
\begin{equation}\label{eq:Hg}
  h_g(P,T) = h_o(P,T) + L + c_{p,g}\,(T - T_{\mathrm{ref}})
            - c_{p,o}\,(T - T_{\mathrm{ref}})
            + v_g\,(1 - T\,\alpha_g)\,(P - P_{\mathrm{ref}})
            - v_o\,(1 - T\,\alpha_o)\,(P - P_{\mathrm{ref}}).
\end{equation}
Упрощая (вычитаем $h_o$ из обеих частей), получаем разность энтальпий фаз:
\begin{equation}\label{eq:Hg_diff}
  \boxed{\;
    h_g(P,T) - h_o(P,T) = L + (c_{p,g} - c_{p,o})\,(T - T_{\mathrm{ref}})
    + (v_g - v_o)\,(1 - T\,\alpha_{\mathrm{eff}})\,(P - P_{\mathrm{ref}}),
  \;}
\end{equation}
где
\begin{equation}\label{eq:alpha_eff}
  \alpha_{\mathrm{eff}} = \frac{v_g\,\alpha_g - v_o\,\alpha_o}{v_g - v_o}.
\end{equation}

\paragraph{Замечание (упрощение для газа).}
Для идеального газа $\alpha_g = 1/T$, и тогда $1 - T\,\alpha_g = 0$~---
давленческая поправка к энтальпии газа обращается в нуль. Это хорошее
приближение для попутного газа при умеренных давлениях ($P \lesssim 10$~МПа).
В этом случае:
\begin{equation}\label{eq:Hg_ideal}
  h_g(P,T) - h_o(P,T) \approx L + (c_{p,g} - c_{p,o})\,(T - T_{\mathrm{ref}})
    - v_o\,(1 - T\,\alpha_o)\,(P - P_{\mathrm{ref}}).
\end{equation}

%-----------------------------------------------------------------------------
\subsection{Подстановка в энтальпийный баланс}
%-----------------------------------------------------------------------------

Подставляя~\eqref{eq:Hmix} в~\eqref{eq:enthalpy}:
\begin{equation}\label{eq:balance_full}
  (1 - x_{g,1})\,h_o(P_1,T_1) + x_{g,1}\,h_g(P_1,T_1)
  = (1 - x_{g,2})\,h_o(P_2,T_2) + x_{g,2}\,h_g(P_2,T_2),
\end{equation}
где $x_{g,1} = x_g(P_1,T_1)$, $x_{g,2} = x_g(P_2,T_2)$.

Это одно нелинейное уравнение относительно одного неизвестного~$T_2$
(давление~$P_2$ задано). Все величины в правой части~--- известные функции
$P_2$ и~$T_2$:
\begin{itemize}[nosep]
  \item $x_{g,2}$ вычисляется по~\eqref{eq:xg} через $R_s(P_2, T_2)$;
  \item $h_o(P_2,T_2)$ вычисляется по~\eqref{eq:Ho};
  \item $h_g(P_2,T_2)$ вычисляется по~\eqref{eq:Hg_diff} и~\eqref{eq:Ho}.
\end{itemize}

%-----------------------------------------------------------------------------
\subsection{Замыкающие корреляции}
%-----------------------------------------------------------------------------

Модель замыкается эмпирическими корреляциями для $R_s$, $\rho_o$, $\rho_g$.

\paragraph{Газосодержание (Стендинг).}
\begin{equation}\label{eq:standing_full}
  R_s(P,T) = \gamma_g\,\left[\frac{P}{a}\,
    10^{\,b\,\gamma_{\mathrm{API}} - c\,T_F}\right]^d,
\end{equation}
где $T_F$~--- температура в~${}^\circ\mathrm{F}$; константы:
$a = 18{,}2$, $b = 0{,}0125$, $c = 0{,}00091$, $d = 1{,}2048$;
давление $P$ в~psia.

\paragraph{Плотность нефти (Стендинг--Кац).}
\begin{equation}\label{eq:rho_o}
  \rho_o(P,T) = \frac{\rho_{o,\text{ст}} + \rho_{g,\text{ст}}\,R_s(P,T)}{B_o(P,T)},
\end{equation}
где $B_o$~--- объёмный коэффициент нефти.

\paragraph{Плотность газа.}
При умеренных давлениях (до ${\sim}20$~МПа) можно использовать закон
реального газа:
\begin{equation}\label{eq:rho_g}
  \rho_g(P,T) = \frac{P\,M_g}{Z(P,T)\,R\,T},
\end{equation}
где $M_g$~--- молярная масса газа, $Z$~--- коэффициент сверхсжимаемости
(определяется по корреляции Стендинга--Кац или Дранчука--Абу-Кассема),
$R = 8{,}314$~Дж/(моль$\cdot$К).

%-----------------------------------------------------------------------------
\subsection{Алгоритм решения}
%-----------------------------------------------------------------------------

\begin{enumerate}
  \item \textbf{Входные данные:} $P_1$, $T_1$, $P_2$; свойства нефти
        ($\gamma_{\mathrm{API}}$, $\rho_{o,\text{ст}}$, $c_{p,o}$,
        $\alpha_o$); свойства газа ($\gamma_g$, $M_g$, $c_{p,g}$); $L$.

  \item \textbf{Вычислить состояние в сечении~1:}
    \begin{enumerate}[nosep]
      \item $R_{s,1} = R_s(P_1, T_1)$ по~\eqref{eq:standing_full};
      \item $x_{g,1}$ по~\eqref{eq:xg};
      \item $h_o(P_1, T_1)$ по~\eqref{eq:Ho};
      \item $h_g(P_1, T_1) = h_o(P_1,T_1) + [h_g - h_o]$ по~\eqref{eq:Hg_diff};
      \item $h_{\mathrm{mix},1}$ по~\eqref{eq:Hmix}.
    \end{enumerate}

  \item \textbf{Итерационный поиск $T_2$ методом бисекции:}
    \begin{enumerate}[nosep]
      \item Задать начальный интервал $[T_{\min}, T_{\max}]$, например
            $[T_1 - 50, \; T_1]$;
      \item Для пробного $T_2^{(k)}$:
        \begin{itemize}[nosep]
          \item вычислить $R_s(P_2, T_2^{(k)})$, $x_{g,2}^{(k)}$,
                $\rho_o^{(k)}$, $\rho_g^{(k)}$;
          \item вычислить $h_o(P_2, T_2^{(k)})$,
                $h_g(P_2, T_2^{(k)})$;
          \item вычислить $h_{\mathrm{mix},2}^{(k)}$;
        \end{itemize}
      \item Вычислить невязку
            $\varepsilon^{(k)} = h_{\mathrm{mix},2}^{(k)} - h_{\mathrm{mix},1}$;
      \item Обновить интервал бисекции по знаку $\varepsilon^{(k)}$;
      \item Повторять, пока $|\varepsilon^{(k)}| > \delta$ (например,
            $\delta = 10$~Дж/кг, что соответствует ${\sim}0{,}005\;{}^\circ$С).
    \end{enumerate}

  \item \textbf{Результат:} $T_2$, а также $x_{g,2}$, $R_{s,2}$~--- полное
        термодинамическое состояние в сечении~2.
\end{enumerate}

\paragraph{Сходимость.}
Функция $h_{\mathrm{mix}}(P_2, T)$ монотонно возрастает по~$T$ (так как
$c_{p} > 0$), поэтому бисекция гарантированно сходится. На практике
достаточно 15--20~итераций для точности $10^{-4}\;{}^\circ$С.

%=============================================================================
\section{Модель~4. Нестационарный энтальпийный баланс}
%=============================================================================

\subsection{Постановка задачи}

Рассматривается контрольный объём~$V$ (например, призабойная зона
скважины), ограниченный двумя контурами (рис.~\ref{fig:CV}):
\begin{itemize}[nosep]
  \item \textbf{внешний контур~$\Gamma_2$}~--- через него поступает живая
        нефть из пласта при давлении~$P_2(t)$ и температуре~$T_2(t)$
        (пластовые условия);
  \item \textbf{внутренний контур~$\Gamma_1$}~--- через него отбирается
        продукция (дегазированная нефть и свободный газ) при забойном
        давлении~$P_1(t)$ и температуре~$T_1(t)$.
\end{itemize}

Давление в объёме меняется во времени за счёт отбора. Когда давление
падает ниже давления насыщения, в объёме начинается разгазирование.
Требуется определить эволюцию температуры $T(t)$ в контрольном объёме.

\begin{figure}[h]
\centering
\setlength{\unitlength}{1mm}
\begin{picture}(100,50)
  % Внешний контур
  \put(10,5){\framebox(80,40){}}
  \put(50,25){\makebox(0,0){$V,\; P(t),\; T(t)$}}
  \put(50,25){\makebox(0,0)[b]{\rule{0pt}{8mm}$m_o(t),\; m_g(t)$}}
  % Внешний контур — подпись
  \put(90,25){\vector(1,0){0}}
  \put(100,25){\makebox(0,0)[l]{\small $\Gamma_2$}}
  % Приток
  \put(90,35){\vector(-1,0){12}}
  \put(97,37){\makebox(0,0)[b]{\small $\dot{m}_{\text{in}}$}}
  % Внутренний контур — подпись
  \put(0,25){\makebox(0,0)[r]{\small $\Gamma_1$}}
  % Отток
  \put(10,35){\vector(-1,0){12}}
  \put(3,37){\makebox(0,0)[b]{\small $\dot{m}_{\text{out}}$}}
  % Подписи условий
  \put(97,30){\makebox(0,0)[l]{\scriptsize $P_2, T_2$}}
  \put(3,30){\makebox(0,0)[r]{\scriptsize $P_1, T_1$}}
\end{picture}
\caption{Контрольный объём с внешним ($\Gamma_2$) и внутренним ($\Gamma_1$)
контурами.}
\label{fig:CV}
\end{figure}

\paragraph{Допущения.}
\begin{enumerate}[nosep]
  \item Контрольный объём термодинамически однородный (<<хорошо
        перемешан>>): давление и температура в нём одинаковы
        и равны $P(t)$, $T(t)$.
  \item Фазы в объёме находятся в термодинамическом равновесии.
  \item Теплообмен с вмещающими породами пренебрежимо мал (адиабатический
        вариант) либо задаётся слагаемым~$\dot{Q}$.
  \item Объём пустотного пространства $V_p$ постоянен.
\end{enumerate}

%-----------------------------------------------------------------------------
\subsection{Дополнительные обозначения}
%-----------------------------------------------------------------------------

\begin{center}
\renewcommand{\arraystretch}{1.25}
\begin{tabular}{cl}
  \hline
  Символ & Смысл \\
  \hline
  $t$                   & время, с \\
  $V_p$                 & объём пустотного пространства (пор), м$^3$ \\
  $P(t)$, $T(t)$        & давление и температура в контрольном объёме \\
  $P_2(t)$, $T_2(t)$    & давление и температура на внешнем контуре \\
  $P_1(t)$, $T_1(t)$    & давление и температура на внутреннем контуре (забой) \\
  $\dot{m}_{\text{in}}$ & массовый расход притока через $\Gamma_2$, кг/с \\
  $\dot{m}_{\text{out}}$& массовый расход отбора через $\Gamma_1$, кг/с \\
  $S_o$, $S_g$          & насыщенности (объёмные доли) нефтяной и газовой фаз \\
  $\dot{Q}$             & мощность теплообмена с окружающими породами, Вт \\
  $u_o$, $u_g$          & удельные внутренние энергии фаз, Дж/кг \\
  \hline
\end{tabular}
\end{center}

%-----------------------------------------------------------------------------
\subsection{Баланс массы}
%-----------------------------------------------------------------------------

Полная масса флюида в контрольном объёме:
\begin{equation}\label{eq:M_CV}
  M(t) = V_p\,\bigl[S_o(t)\,\rho_o(t) + S_g(t)\,\rho_g(t)\bigr],
  \qquad S_o + S_g = 1.
\end{equation}

Через внешний контур входит живая нефть (однофазная, при $P_2 > P_{\text{нас}}$)
с массовым расходом~$\dot{m}_{\text{in}}(t)$.
Через внутренний контур выходит двухфазная смесь с расходом~$\dot{m}_{\text{out}}(t)$.

Баланс полной массы:
\begin{equation}\label{eq:mass_total}
  \boxed{\;\frac{dM}{dt} = \dot{m}_{\text{in}}(t) - \dot{m}_{\text{out}}(t).\;}
\end{equation}

Кроме полной массы, удобно записать баланс массы лёгкого компонента
(газа, считая и растворённый, и свободный). Масса лёгкого компонента
в объёме:
\begin{equation}\label{eq:Mg_CV}
  M_{\text{л}}(t) = V_p\,S_g\,\rho_g
    + V_p\,S_o\,\rho_o\,\frac{\rho_{g,\text{ст}}\,R_s(P,T)}
      {\rho_{o,\text{ст}} + \rho_{g,\text{ст}}\,R_s(P,T)}.
\end{equation}
Первое слагаемое~--- масса свободного газа; второе~--- масса газа,
растворённого в нефтяной фазе.

Баланс лёгкого компонента:
\begin{equation}\label{eq:mass_light}
  \frac{dM_{\text{л}}}{dt}
    = \dot{m}_{\text{in}}\,x_{g,\text{in}}
    - \dot{m}_{\text{out}}\,x_{g,\text{out}},
\end{equation}
где $x_{g,\text{in}}$~--- массовая доля лёгкого компонента в притоке
(определяется по $R_s(P_2, T_2)$), $x_{g,\text{out}}$~--- массовая доля
лёгкого компонента в отбираемой продукции (определяется по $R_s(P_1, T_1)$
и доле свободного газа).

%-----------------------------------------------------------------------------
\subsection{Баланс энергии}
%-----------------------------------------------------------------------------

Внутренняя энергия флюида в контрольном объёме:
\begin{equation}\label{eq:U_CV}
  U(t) = V_p\,\bigl[S_o\,\rho_o\,u_o(P,T) + S_g\,\rho_g\,u_g(P,T)\bigr].
\end{equation}

Первое начало термодинамики для открытой системы:
\begin{equation}\label{eq:energy}
  \boxed{\;
    \frac{dU}{dt}
    = \dot{m}_{\text{in}}\,h_{\text{in}}
    - \dot{m}_{\text{out}}\,h_{\text{out}}
    + \dot{Q},
  \;}
\end{equation}
где
\begin{itemize}[nosep]
  \item $h_{\text{in}} = h_{\mathrm{mix}}(P_2, T_2)$~--- удельная энтальпия
        притока (живая нефть при пластовых условиях; вычисляется
        по~\eqref{eq:Hmix}, \eqref{eq:Ho}, \eqref{eq:Hg_diff});
  \item $h_{\text{out}} = h_{\mathrm{mix}}(P_1, T_1)$~--- удельная энтальпия
        отбираемой продукции;
  \item $\dot{Q}$~--- теплоприток от вмещающих пород ($\dot{Q} > 0$~--- нагрев,
        $\dot{Q} = 0$~--- адиабатический случай).
\end{itemize}

\paragraph{Связь $u$ и $h$.}
Удельная внутренняя энергия связана с энтальпией:
\begin{equation}\label{eq:u_from_h}
  u = h - P\,v = h - \frac{P}{\rho}.
\end{equation}
Поэтому \eqref{eq:U_CV} можно выразить через энтальпии:
\begin{equation}\label{eq:U_via_h}
  U = V_p\,\Bigl[S_o\,\rho_o\,\Bigl(h_o - \frac{P}{\rho_o}\Bigr)
    + S_g\,\rho_g\,\Bigl(h_g - \frac{P}{\rho_g}\Bigr)\Bigr]
    = V_p\,\bigl[S_o\,\rho_o\,h_o + S_g\,\rho_g\,h_g - P\bigr].
\end{equation}

%-----------------------------------------------------------------------------
\subsection{Замыкание системы}
%-----------------------------------------------------------------------------

Система из трёх уравнений~\eqref{eq:mass_total},
\eqref{eq:mass_light}, \eqref{eq:energy} содержит три неизвестных:
\begin{equation}\label{eq:unknowns}
  P(t), \quad T(t), \quad S_o(t) \quad (\text{или } S_g = 1 - S_o).
\end{equation}

Все остальные величины~--- функции $(P, T, S_o)$:
\begin{itemize}[nosep]
  \item $R_s(P,T)$~--- по корреляции Стендинга~\eqref{eq:standing_full};
  \item $\rho_o(P,T)$~--- по~\eqref{eq:rho_o};\quad
        $\rho_g(P,T)$~--- по~\eqref{eq:rho_g};
  \item $h_o(P,T)$~--- по~\eqref{eq:Ho};\quad
        $h_g(P,T)$~--- по~\eqref{eq:Hg_diff} и~\eqref{eq:Ho};
  \item $x_g(P,T)$~--- по~\eqref{eq:xg}.
\end{itemize}

Задаваемые внешние данные (история разработки):
\begin{itemize}[nosep]
  \item $P_2(t)$, $T_2(t)$~--- условия на внешнем контуре (пластовые);
  \item $\dot{m}_{\text{in}}(t)$~--- расход притока
        (или $P_1(t)$ и формула притока, например закон Дарси);
  \item $\dot{m}_{\text{out}}(t)$~--- расход отбора (или режим работы скважины).
\end{itemize}

%-----------------------------------------------------------------------------
\subsection{Связь с моделью~3 (стационарный предел)}
%-----------------------------------------------------------------------------

Если состояние в объёме не меняется во времени ($dM/dt = 0$, $dU/dt = 0$),
из~\eqref{eq:mass_total} следует $\dot{m}_{\text{in}} =
\dot{m}_{\text{out}} = \dot{m}$, а из~\eqref{eq:energy}:
\begin{equation}
  h_{\mathrm{mix}}(P_2, T_2) = h_{\mathrm{mix}}(P_1, T_1),
\end{equation}
что в точности совпадает с уравнением~\eqref{eq:enthalpy} модели~3.

%-----------------------------------------------------------------------------
\subsection{Численное решение}
%-----------------------------------------------------------------------------

Система ОДУ~\eqref{eq:mass_total}, \eqref{eq:mass_light}, \eqref{eq:energy}
решается методом Эйлера или Рунге--Кутты. На каждом временном шаге~$\Delta t$:

\begin{enumerate}
  \item \textbf{Известно:} $P^n$, $T^n$, $S_o^n$ (состояние на шаге~$n$).

  \item \textbf{Вычислить свойства флюидов} при $(P^n, T^n)$:
        $R_s^n$, $\rho_o^n$, $\rho_g^n$, $h_o^n$, $h_g^n$, $x_g^n$.

  \item \textbf{Вычислить правые части:}
    \begin{align}
      \dot{M}^n &= \dot{m}_{\text{in}}^n - \dot{m}_{\text{out}}^n,
        \label{eq:rhs_M} \\
      \dot{M}_{\text{л}}^n &= \dot{m}_{\text{in}}^n\,x_{g,\text{in}}^n
        - \dot{m}_{\text{out}}^n\,x_{g,\text{out}}^n,
        \label{eq:rhs_Ml} \\
      \dot{U}^n &= \dot{m}_{\text{in}}^n\,h_{\text{in}}^n
        - \dot{m}_{\text{out}}^n\,h_{\text{out}}^n + \dot{Q}^n.
        \label{eq:rhs_U}
    \end{align}

  \item \textbf{Обновить интегральные величины:}
    \begin{equation}\label{eq:euler}
      M^{n+1} = M^n + \dot{M}^n\,\Delta t, \quad
      M_{\text{л}}^{n+1} = M_{\text{л}}^n + \dot{M}_{\text{л}}^n\,\Delta t, \quad
      U^{n+1} = U^n + \dot{U}^n\,\Delta t.
    \end{equation}

  \item \textbf{Восстановить} $P^{n+1}$, $T^{n+1}$, $S_o^{n+1}$ из
        системы трёх алгебраических уравнений:
    \begin{align}
      V_p\,\bigl[S_o\,\rho_o(P,T) + (1-S_o)\,\rho_g(P,T)\bigr] &= M^{n+1},
        \label{eq:inv_M} \\
      M_{\text{л}}(P,T,S_o) &= M_{\text{л}}^{n+1},
        \label{eq:inv_Ml} \\
      V_p\,\bigl[S_o\,\rho_o\,h_o + (1-S_o)\,\rho_g\,h_g - P\bigr] &= U^{n+1}.
        \label{eq:inv_U}
    \end{align}
        Эта система решается итерационно (метод Ньютона, 3--5~итераций).
\end{enumerate}

\paragraph{Частный случай: заданное $P(t)$.}
Если история давления $P(t)$ известна (например, из данных манометрии),
система упрощается: остаются два неизвестных $T(t)$ и $S_o(t)$,
и шаг~5 сводится к системе двух уравнений~\eqref{eq:inv_Ml}
и~\eqref{eq:inv_U}, что существенно ускоряет решение.

%=============================================================================
\section{Модель~5. Пространственно-распределённая задача\\
с упрощённой гидродинамикой}
%=============================================================================

\subsection{Идея упрощения}

Полная двухфазная модель фильтрации (закон Дарси для каждой фазы,
ОФП, два уравнения неразрывности) существенно усложняет задачу и
требует данных, которые редко известны с хорошей точностью (кривые
ОФП, вязкости двух фаз и т.\,д.). Между тем основная физика
охлаждения --- в термодинамике фазового перехода, а не в деталях
двухфазного течения.

Ключевое упрощение: \textbf{считать поле давления $P(r,t)$ заданным},
а решать только уравнение энергии. Давление можно получить одним из
способов:
\begin{enumerate}[nosep]
  \item из данных манометрии (интерполяция замеров);
  \item из аналитического решения однофазной задачи фильтрации;
  \item из простой численной модели однофазного притока.
\end{enumerate}

При заданном $P(r,t)$ газосодержание $R_s$, плотности, энтальпии~---
всё определяется локально, и задача сводится к \textbf{одному PDE}
для температуры~$T(r,t)$.

\subsection{Геометрия и обозначения}

Радиально-симметричная постановка: $r_w \leqslant r \leqslant r_e$,
$r_w$~--- радиус скважины, $r_e$~--- радиус контура питания.

\begin{center}
\renewcommand{\arraystretch}{1.25}
\begin{tabular}{cl}
  \hline
  Символ & Смысл \\
  \hline
  $r$                   & радиальная координата, м \\
  $r_w$, $r_e$          & радиус скважины и контура питания \\
  $\phi$                & пористость \\
  $k$                   & абсолютная проницаемость, м$^2$ \\
  $\mu$                 & вязкость пластовой нефти (однофазной), Па$\cdot$с \\
  $v(r,t)$              & скорость фильтрации суммарного потока, м/с \\
  $\lambda_{\mathrm{eff}}$ & эффективная теплопроводность
                             (флюид + порода), Вт/(м$\cdot$К) \\
  $(\rho c_p)_s$        & объёмная теплоёмкость скелета породы, Дж/(м$^3\cdot$К) \\
  $Q_{\text{дег}}$      & объёмная мощность источника холода
                           (разгазирование), Вт/м$^3$ \\
  \hline
\end{tabular}
\end{center}

%-----------------------------------------------------------------------------
\subsection{Поле давления (задано)}
%-----------------------------------------------------------------------------

Давление $P(r,t)$ считается известным. Простейший вариант~---
однофазная радиальная фильтрация несжимаемой жидкости:
\begin{equation}\label{eq:P_ss}
  P(r) = P_{\text{заб}} + \frac{Q\,\mu}{2\pi\,k\,h_{\text{пл}}}\,
         \ln\frac{r}{r_w},
\end{equation}
где $Q$~--- объёмный дебит, $h_{\text{пл}}$~--- толщина пласта.
Для нестационарного случая~--- решение Тейса или данные манометрии.

Скорость фильтрации суммарного потока определяется из баланса массы:
\begin{equation}\label{eq:v_total}
  v(r,t) = \frac{Q(t)}{2\pi\,r\,h_{\text{пл}}}.
\end{equation}
Формула~\eqref{eq:v_total} точна для несжимаемого потока и является
хорошим приближением при умеренной сжимаемости.

%-----------------------------------------------------------------------------
\subsection{Локальное термодинамическое состояние}
%-----------------------------------------------------------------------------

В каждой точке $(r,t)$ при известном давлении $P(r,t)$ и искомой
температуре $T(r,t)$ определяются:

\paragraph{Газосодержание и массовая доля газа.}
\begin{equation}\label{eq:Rs_local}
  R_s = R_s\bigl(P(r,t),\,T(r,t)\bigr), \qquad
  x_g = \frac{\rho_{g,\text{ст}}\,\bigl[R_{s,\max} - R_s\bigr]}
             {\rho_{o,\text{ст}} + \rho_{g,\text{ст}}\,R_{s,\max}}.
\end{equation}

\paragraph{Плотности фаз.}
$\rho_o(P,T)$ по~\eqref{eq:rho_o}, $\rho_g(P,T)$ по~\eqref{eq:rho_g}.

\paragraph{Средняя плотность смеси.}
\begin{equation}\label{eq:rho_mix}
  \rho_{\mathrm{mix}} = S_o\,\rho_o + S_g\,\rho_g.
\end{equation}

\paragraph{Насыщенность газом.}
При известных $x_g$, $\rho_o$, $\rho_g$ насыщенность определяется
алгебраически (без решения уравнения фильтрации):
\begin{equation}\label{eq:Sg_from_xg}
  S_g = \frac{x_g\,\rho_o}{x_g\,\rho_o + (1-x_g)\,\rho_g}, \qquad
  S_o = 1 - S_g.
\end{equation}
Это соотношение следует из равновесия фаз в предположении, что свободный
газ не мигрирует относительно нефти (нулевое относительное проскальзывание
фаз~--- допустимо при малых $S_g$, характерных для начала разгазирования).

\paragraph{Энтальпия смеси.}
\begin{equation}\label{eq:hmix_local}
  h_{\mathrm{mix}} = (1 - x_g)\,h_o + x_g\,h_g,
\end{equation}
где $h_o$, $h_g$~--- по формулам~\eqref{eq:Ho}, \eqref{eq:Hg_diff}.

%-----------------------------------------------------------------------------
\subsection{Источник холода: энергия разгазирования}
%-----------------------------------------------------------------------------

При снижении давления в точке $(r,t)$ газ выделяется из нефти.
Объёмная мощность отбора тепла на разгазирование:
\begin{equation}\label{eq:Q_degas}
  \boxed{\;
    Q_{\text{дег}} = -\,\phi\,\rho_{\mathrm{mix}}\,
    \frac{\partial h_{\mathrm{mix}}}{\partial P}\bigg|_{T}\,
    \frac{\partial P}{\partial t}.
  \;}
\end{equation}
Здесь $\partial h_{\mathrm{mix}}/\partial P\big|_T$ включает как
эффект Джоуля--Томсона (изменение энтальпии фаз с давлением), так и
скрытую теплоту дегазации (через зависимость $x_g$ от~$P$):
\begin{equation}\label{eq:dh_dP}
  \frac{\partial h_{\mathrm{mix}}}{\partial P}\bigg|_{T}
  = (1-x_g)\,\frac{\partial h_o}{\partial P}\bigg|_T
  + x_g\,\frac{\partial h_g}{\partial P}\bigg|_T
  + (h_g - h_o)\,\frac{\partial x_g}{\partial P}\bigg|_T.
\end{equation}
Последнее слагаемое~--- главный вклад при $P < P_{\text{нас}}$:
оно описывает затраты энергии на перевод массы из жидкой фазы в
газовую ($\partial x_g/\partial P < 0$, $h_g - h_o > 0$, поэтому вклад
отрицательный~--- охлаждение).

%-----------------------------------------------------------------------------
\subsection{Уравнение энергии}
%-----------------------------------------------------------------------------

Единственное PDE модели:
\begin{equation}\label{eq:energy_pde}
  \boxed{\;
    \underbrace{(\rho c_p)_{\mathrm{eff}}\,
      \frac{\partial T}{\partial t}}_{\text{накопление}}
    + \underbrace{\rho_{\mathrm{mix}}\,c_{p,\mathrm{mix}}\,v\,
      \frac{\partial T}{\partial r}}_{\text{конвекция}}
    = \underbrace{\frac{1}{r}\,\frac{\partial}{\partial r}
      \Bigl(r\,\lambda_{\mathrm{eff}}\,
      \frac{\partial T}{\partial r}\Bigr)}_{\text{теплопроводность}}
    + \underbrace{Q_{\text{дег}}}_{\text{разгазирование}}
    + \underbrace{Q_{\mathrm{JT}}}_{\text{Джоуль--Томсон}},
  \;}
\end{equation}
где:
\begin{itemize}[nosep]
  \item эффективная объёмная теплоёмкость среды (флюид + порода):
    \begin{equation}\label{eq:rhoCp_eff}
      (\rho c_p)_{\mathrm{eff}} = \phi\,\rho_{\mathrm{mix}}\,
        c_{p,\mathrm{mix}} + (1-\phi)\,(\rho c_p)_s;
    \end{equation}
  \item теплоёмкость смеси:
    \begin{equation}\label{eq:cp_mix}
      c_{p,\mathrm{mix}} = (1 - x_g)\,c_{p,o} + x_g\,c_{p,g};
    \end{equation}
  \item источник Джоуля--Томсона (охлаждение при адиабатическом
        расширении однофазной нефти):
    \begin{equation}\label{eq:Q_JT}
      Q_{\mathrm{JT}} = -\,\phi\,\rho_{\mathrm{mix}}\,
        c_{p,\mathrm{mix}}\,\mu_{\mathrm{JT}}\,
        \frac{\partial P}{\partial t};
    \end{equation}
  \item источник разгазирования $Q_{\text{дег}}$~---
        по~\eqref{eq:Q_degas}.
\end{itemize}

\paragraph{Структура уравнения.}
Уравнение~\eqref{eq:energy_pde}~--- линейное параболическое PDE
относительно~$T(r,t)$ с коэффициентами, зависящими от~$P(r,t)$
(которое задано). Все <<тяжёлые>> эффекты разгазирования входят
через источниковые члены $Q_{\text{дег}}$ и $Q_{\mathrm{JT}}$,
которые вычисляются по заданному полю давления.

%-----------------------------------------------------------------------------
\subsection{Начальные и граничные условия}
%-----------------------------------------------------------------------------

\paragraph{Начальные условия} ($t = 0$):
\begin{equation}\label{eq:IC}
  T(r, 0) = T_{\text{пл}} \quad \text{(пластовая температура)}.
\end{equation}

\paragraph{Внешний контур} ($r = r_e$):
\begin{equation}\label{eq:BC_re}
  T(r_e, t) = T_{\text{пл}} \quad \text{(невозмущённый пласт)}.
\end{equation}

\paragraph{Стенка скважины} ($r = r_w$):
\begin{equation}\label{eq:BC_rw}
  \frac{\partial T}{\partial r}\bigg|_{r=r_w} = 0 \quad
  \text{(адиабатическое условие)}
\end{equation}
или заданная температура/теплопоток.

%-----------------------------------------------------------------------------
\subsection{Связь с предыдущими моделями}
%-----------------------------------------------------------------------------

\paragraph{Предел $\to$ модель~3.}
В стационарном случае ($\partial T/\partial t = 0$,
$\partial P/\partial t = 0$) и при пренебрежении теплопроводностью
уравнение~\eqref{eq:energy_pde} вырождается в
\begin{equation}
  v\,\frac{d\,h_{\mathrm{mix}}}{dr} = 0
  \quad\Longrightarrow\quad
  h_{\mathrm{mix}} = \mathrm{const},
\end{equation}
что совпадает с энтальпийным балансом~\eqref{eq:enthalpy} модели~3.

\paragraph{Предел $\to$ модель~4.}
Интегрирование~\eqref{eq:energy_pde} по объёму
$\int_{r_w}^{r_e} (\cdot)\,2\pi r\,h_{\text{пл}}\,dr$
с заменой пространственно распределённых величин их средними
воспроизводит энергетический баланс~\eqref{eq:energy} модели~4.

%-----------------------------------------------------------------------------
\subsection{Численное решение}
%-----------------------------------------------------------------------------

Уравнение~\eqref{eq:energy_pde} решается стандартными методами
для параболических PDE.

\paragraph{Сетка.}
Логарифмическая по~$r$ (сгущение к скважине):
\begin{equation}\label{eq:grid}
  r_j = r_w\,\left(\frac{r_e}{r_w}\right)^{j/N}, \quad j = 0, 1, \ldots, N.
\end{equation}

\paragraph{Дискретизация по времени.}
Неявная схема (безусловно устойчива):
\begin{equation}\label{eq:implicit}
  (\rho c_p)_{\mathrm{eff},j}\,\frac{T_j^{n+1} - T_j^n}{\Delta t}
  = \mathcal{L}[T^{n+1}]_j + Q_{\text{дег},j}^n + Q_{\mathrm{JT},j}^n,
\end{equation}
где $\mathcal{L}$~--- дискретный оператор конвекции и диффузии.
На каждом шаге решается \textbf{трёхдиагональная} система $N$~уравнений
(метод прогонки, $O(N)$ операций).

\paragraph{Оценка сложности.}
При $N = 50\text{--}100$ узлов и ${\sim}10^3$ шагов по времени задача
решается за доли секунды. Это на порядки быстрее полной двухфазной
модели, где на каждом шаге требуется решать нелинейную систему~$3N$
уравнений.

%=============================================================================
\section{Модель~6. GORM с учётом дегазации и температуры}
%=============================================================================

\subsection{Исходная модель GORM}

Модель GORM (Gas/Oil Ratio Model), предложенная
Konieczek~\cite{konieczek1990} и развитая
Mjaavatten et~al.~\cite{mjaavatten2006}, описывает эволюцию
газонефтяного контакта (ГНК) вблизи горизонтальной скважины в
нефтяной оторочке. Ключевое упрощение: поток нефти к скважине
является горизонтальным и приводится в движение только
гравитацией (разностью плотностей газа и нефти).

Единственная неизвестная~--- толщина нефтяного слоя $h(x,t)$,
удовлетворяющая уравнению Дюпюи--Форхгеймера:
\begin{equation}\label{eq:gorm_base}
  \boxed{\;
    \frac{\partial h}{\partial t}
    = \alpha\,\frac{\partial}{\partial x}
      \left(h\,\frac{\partial h}{\partial x}\right),
    \qquad
    \alpha = \frac{k\,\Delta\rho\,g}{\mu_o\,\varphi},
  \;}
\end{equation}
где $x$~--- горизонтальная координата (перпендикулярно скважине),
$k$~--- горизонтальная проницаемость,
$\Delta\rho = \rho_o - \rho_g$~--- разность плотностей,
$\mu_o$~--- вязкость нефти, $\varphi$~--- эффективная пористость.

\paragraph{Допущения исходной GORM.}
\begin{enumerate}[nosep]
  \item Давление в газовой шапке постоянно.
  \item Флюиды несжимаемы; свойства ($\rho_o$, $\mu_o$) постоянны.
  \item Поток чисто горизонтальный (приближение Дюпюи).
  \item Капиллярные эффекты отсутствуют; ГНК~--- резкая граница.
  \item Водонефтяной контакт неподвижен.
\end{enumerate}

Модель замыкается граничным условием на скважине (дебит нефти
пропорционален $h\,\partial h/\partial x$) и условием непротекания
на внешнем контуре.

%-----------------------------------------------------------------------------
\subsection{Обозначения модели~6}
%-----------------------------------------------------------------------------

\begin{center}
\renewcommand{\arraystretch}{1.25}
\begin{tabular}{cl}
  \hline
  Символ & Смысл \\
  \hline
  $x$                    & горизонтальная координата $\perp$ скважине, м \\
  $y$                    & координата вдоль скважины, м \\
  $h(x,y,t)$            & толщина нефтяного слоя (высота ГНК), м \\
  $h_0$                  & начальная толщина нефтяного слоя, м \\
  $W$                    & полуширина области дренирования, м \\
  $L$                    & длина горизонтального участка скважины, м \\
  $\alpha$               & коэффициент <<гравитационной диффузии>>, м$^2$/с \\
  $\varphi$              & эффективная пористость (подвижная нефть) \\
  $\Delta\rho$           & $\rho_o - \rho_g$, кг/м$^3$ \\
  $\delta(x,y,t)$        & безразмерный уровень нефти у скважины: $(z_w - h)/d_w$ \\
  $\gamma$               & отношение индексов продуктивности газа и нефти: $J_g/J_o$ \\
  $B_o(P)$               & объёмный коэффициент нефти \\
  $B_g(P)$               & объёмный коэффициент газа \\
  $T(x,t)$               & температура, К \\
  \hline
\end{tabular}
\end{center}

%-----------------------------------------------------------------------------
\subsection{Обобщение: учёт дегазации}
%-----------------------------------------------------------------------------

При снижении пластового давления ниже давления насыщения
$P_{\text{нас}}$ газ выделяется не только из газовой шапки, но и
\textit{внутри} нефтяного слоя. Это приводит к трём эффектам:
\begin{enumerate}[nosep]
  \item нефть <<усыхает>> (объёмный коэффициент $B_o$ уменьшается);
  \item в нефтяном слое появляется свободный газ
        с насыщенностью $S_g(P,T)$;
  \item подвижность нефти снижается из-за присутствия газовых пузырьков.
\end{enumerate}

\paragraph{Давление в нефтяном слое.}
В рамках приближения Дюпюи (вертикальное равновесие) давление
определяется гидростатически:
\begin{equation}\label{eq:gorm_pressure}
  P(x,z,t) = P_{\text{газ}} + \rho_g\,g\,(h_0 - h)
              + \rho_o\,g\,(h - z),
\end{equation}
где $P_{\text{газ}}$~--- давление в газовой шапке (может медленно
снижаться при разработке), $z$~--- вертикальная координата,
отсчитанная вверх от подошвы.

Среднее давление в нефтяном столбе по вертикали:
\begin{equation}\label{eq:gorm_Pavg}
  \bar{P}(x,t) = P_{\text{газ}} + \rho_g\,g\,(h_0 - h)
                + \tfrac{1}{2}\,\rho_o\,g\,h.
\end{equation}

\paragraph{Локальное газосодержание.}
При $\bar{P} < P_{\text{нас}}$ газ выделяется. Газосодержание:
\begin{equation}\label{eq:gorm_Rs}
  R_s\bigl(\bar{P},T\bigr) \quad\text{--- по корреляции Стендинга
  \eqref{eq:standing_full}}.
\end{equation}
Насыщенность свободным газом в нефтяном слое:
\begin{equation}\label{eq:gorm_Sg}
  S_g(x,t) = \frac{B_o\,\bigl[R_{s,\max} - R_s(\bar{P},T)\bigr]}
                   {B_g + B_o\,\bigl[R_{s,\max} - R_s(\bar{P},T)\bigr]},
  \qquad S_o = 1 - S_g.
\end{equation}

%-----------------------------------------------------------------------------
\subsection{Модифицированное уравнение для~$h$}
%-----------------------------------------------------------------------------

В исходной GORM объём подвижной нефти на единицу площади
$V = \varphi\,h$. При дегазации нужно учесть, что часть порового
пространства занята свободным газом:
\begin{equation}\label{eq:gorm_V}
  V_o(x,t) = \varphi\,h\,S_o(x,t).
\end{equation}

Обобщённое уравнение Дюпюи--Форхгеймера:
\begin{equation}\label{eq:gorm_modified}
  \boxed{\;
    \frac{\partial}{\partial t}\bigl(\varphi\,h\,S_o\bigr)
    = \frac{\partial}{\partial x}
      \left(\alpha_{\mathrm{eff}}\,h\,
      \frac{\partial h}{\partial x}\right),
  \;}
\end{equation}
где эффективный коэффициент гравитационной диффузии:
\begin{equation}\label{eq:alpha_eff_gorm}
  \alpha_{\mathrm{eff}}(x,t)
    = \frac{k\,k_{r,o}(S_o)\,\Delta\rho\,g}{\mu_o(\bar{P},T)}.
\end{equation}
Здесь $k_{r,o}(S_o)$~--- относительная проницаемость нефти. В простейшем
приближении можно положить $k_{r,o} \approx S_o$ (линейная модель).

\paragraph{Предел $\to$ исходная GORM.}
При $\bar{P} > P_{\text{нас}}$ имеем $S_g = 0$, $S_o = 1$,
$k_{r,o} = 1$, и уравнение~\eqref{eq:gorm_modified} вырождается
в~\eqref{eq:gorm_base}.

%-----------------------------------------------------------------------------
\subsection{Уравнение энергии}
%-----------------------------------------------------------------------------

Температура в нефтяном слое описывается уравнением теплопереноса,
аналогичным~\eqref{eq:energy_pde} модели~5, но записанным в
координатах GORM. Скорость фильтрации берётся из закона Дюпюи:
\begin{equation}\label{eq:gorm_velocity}
  v(x,t) = -\,\frac{k\,k_{r,o}\,\Delta\rho\,g}{\mu_o}\,
            \frac{\partial h}{\partial x}
          = -\,\frac{\alpha_{\mathrm{eff}}}{\varphi}\,
            \frac{\partial h}{\partial x}.
\end{equation}

Уравнение энергии для средней по вертикали температуры $T(x,t)$:
\begin{equation}\label{eq:gorm_energy}
  \boxed{\;
    (\rho c_p)_{\mathrm{eff}}\,\frac{\partial T}{\partial t}
    + \rho_{\mathrm{mix}}\,c_{p,\mathrm{mix}}\,v\,
      \frac{\partial T}{\partial x}
    = \frac{\partial}{\partial x}
      \left(\lambda_{\mathrm{eff}}\,
      \frac{\partial T}{\partial x}\right)
    + Q_{\text{дег}} + Q_{\mathrm{JT}},
  \;}
\end{equation}
где все величины определяются так же, как в модели~5
(уравнения~\eqref{eq:rhoCp_eff}--\eqref{eq:Q_JT}), но давление
берётся из гидростатики~\eqref{eq:gorm_Pavg}, а скорость~---
из~\eqref{eq:gorm_velocity}.

Источниковые члены:
\begin{align}
  Q_{\text{дег}} &= -\,\phi\,\rho_{\mathrm{mix}}\,
    \frac{\partial h_{\mathrm{mix}}}{\partial P}\bigg|_T\,
    \frac{\partial \bar{P}}{\partial t},
    \label{eq:gorm_Qdeg} \\[4pt]
  Q_{\mathrm{JT}} &= -\,\phi\,\rho_{\mathrm{mix}}\,
    c_{p,\mathrm{mix}}\,\mu_{\mathrm{JT}}\,
    \frac{\partial \bar{P}}{\partial t}.
    \label{eq:gorm_QJT}
\end{align}

Производная $\partial\bar{P}/\partial t$ вычисляется
из~\eqref{eq:gorm_Pavg}:
\begin{equation}\label{eq:dPdt_gorm}
  \frac{\partial\bar{P}}{\partial t}
    = \frac{dP_{\text{газ}}}{dt}
    - \bigl(\rho_g - \tfrac{1}{2}\rho_o\bigr)\,g\,
      \frac{\partial h}{\partial t},
\end{equation}
то есть давление в нефтяном слое меняется как из-за снижения давления
в газовой шапке, так и из-за подъёма ГНК (уменьшения~$h$).

%-----------------------------------------------------------------------------
\subsection{Полная система}
%-----------------------------------------------------------------------------

Модель~6 состоит из \textbf{двух связанных PDE}:
\begin{enumerate}[nosep]
  \item уравнение для толщины нефтяного слоя~\eqref{eq:gorm_modified}
        ($h(x,t)$)~--- определяет гидродинамику и давление;
  \item уравнение энергии~\eqref{eq:gorm_energy}
        ($T(x,t)$)~--- определяет температуру.
\end{enumerate}

Связь между ними:
\begin{itemize}[nosep]
  \item $h \to \bar{P}$ по~\eqref{eq:gorm_Pavg}
        $\to R_s, S_g, \rho, h_{\mathrm{mix}}$;
  \item $T \to R_s(\bar{P},T), \mu_o(\bar{P},T)$
        $\to \alpha_{\mathrm{eff}}$.
\end{itemize}

\paragraph{Начальные условия.}
\begin{equation}\label{eq:gorm_IC}
  h(x, 0) = h_0, \qquad T(x, 0) = T_{\text{пл}}.
\end{equation}

\paragraph{Граничные условия.}
На скважине ($x = 0$):
\begin{equation}\label{eq:gorm_BC_well}
  \tilde{q}_o = 2\,\alpha_{\mathrm{eff}}\,\varphi\,h\,
    \frac{\partial h}{\partial x}\bigg|_{x=0}
  \quad\text{(задан дебит нефти)}.
\end{equation}
На внешнем контуре ($x = W$):
\begin{equation}\label{eq:gorm_BC_outer}
  \frac{\partial h}{\partial x}\bigg|_{x=W} = 0
  \quad\text{(непротекание)}, \qquad
  T(W, t) = T_{\text{пл}}
  \quad\text{(невозмущённый пласт)}.
\end{equation}

%-----------------------------------------------------------------------------
\subsection{Газонефтяной фактор (GOR)}
%-----------------------------------------------------------------------------

Как и в исходной GORM, после прорыва газа ($h(0,t) = 0$
на участке скважины) отношение дебитов газа и нефти определяется
через безразмерный уровень $\delta$:
\begin{equation}\label{eq:gorm_GOR_base}
  \frac{\tilde{q}_g}{\tilde{q}_o} = \gamma\,
    \frac{\delta^2}{1 - \delta^2}.
\end{equation}

В обобщённой модели к этому свободному газу из газовой шапки добавляется
газ, выделившийся при дегазации. Полный GOR на поверхности:
\begin{equation}\label{eq:gorm_GOR_total}
  \boxed{\;
    \mathrm{GOR} = \frac{Q_g}{Q_o}
    = \underbrace{R_s(\bar{P},T)}_{\text{растворённый}}
    + \underbrace{\frac{1}{B_o}\,\gamma\,
      \frac{\delta^2}{1-\delta^2}\,B_g}_{\text{газ.\ шапка}}
    + \underbrace{\frac{B_o\,(R_{s,\max} - R_s)}{B_g}
      \cdot\frac{S_g}{S_o}}_{\text{дегазация}}.
  \;}
\end{equation}

%-----------------------------------------------------------------------------
\subsection{Численное решение}
%-----------------------------------------------------------------------------

Оба PDE дискретизируются по~$x$ на равномерной сетке
$x_j = j\,\Delta x$, $j = 0, \ldots, N$.

\paragraph{Уравнение для~$h$.}
Нелинейное параболическое PDE~\eqref{eq:gorm_modified}
решается полунеявной схемой: коэффициент
$\alpha_{\mathrm{eff}}\,h$ берётся с предыдущего шага,
а $h^{n+1}$~--- неявно. Это даёт трёхдиагональную систему.

\paragraph{Уравнение для~$T$.}
Линейное параболическое PDE~\eqref{eq:gorm_energy} при известном~$h$
решается неявной схемой (прогонка), как в модели~5.

\paragraph{Порядок на каждом шаге $\Delta t$.}
\begin{enumerate}[nosep]
  \item Из $h^n$ вычислить $\bar{P}^n$, $R_s^n$, $S_g^n$,
        $\alpha_{\mathrm{eff}}^n$.
  \item Решить~\eqref{eq:gorm_modified} $\to h^{n+1}$
        (трёхдиагональная прогонка).
  \item Вычислить $\partial\bar{P}/\partial t$ из~\eqref{eq:dPdt_gorm},
        источники $Q_{\text{дег}}$, $Q_{\mathrm{JT}}$.
  \item Решить~\eqref{eq:gorm_energy} $\to T^{n+1}$
        (трёхдиагональная прогонка).
\end{enumerate}

На каждом шаге~--- две прогонки, $O(N)$ операций.
При $N = 100$ и $10^3$~шагов по времени решение занимает доли секунды.

%-----------------------------------------------------------------------------
\subsection{Преимущества модели~6}
%-----------------------------------------------------------------------------

\begin{enumerate}[nosep]
  \item \textbf{Самосогласованная гидродинамика}: поле давления и
        скорости вычисляются из PDE для~$h$, а не задаются извне
        (как в модели~5).
  \item \textbf{Минимум входных данных}: те же параметры, что и в
        исходной GORM ($k$, $\varphi$, $\Delta\rho$, $\mu_o$, $W$,
        $h_0$), плюс термодинамические свойства
        ($c_p$, $L$, $\lambda$) и $R_s(P,T)$.
  \item \textbf{Настройка по GOR}: три параметра GORM ($\alpha$,
        $\beta$, $W$) калибруются по истории газового фактора,
        как в~\cite{mjaavatten2006}.
  \item \textbf{Вычислительная эффективность}: два PDE вместо
        трёх ($P$, $T$, $S_o$) в полной модели; $O(N)$ на шаг.
  \item \textbf{Предсказание температуры}: модель позволяет
        прогнозировать температуру на забое и вдоль ствола, что
        полезно для интерпретации данных DTS (Distributed
        Temperature Sensing).
\end{enumerate}

%=============================================================================
\section{Давление и поле скорости в газовой шапке}
%=============================================================================

В модели~6 давление в газовой шапке $P_{\text{газ}}$ до сих пор
считалось заданным (постоянным или медленно снижающимся).
Для замыкания модели и вычисления поля скорости газа в газовой шапке
рассмотрим два подхода разной степени детализации.

\subsection{Дополнительные обозначения}

\begin{center}
\renewcommand{\arraystretch}{1.25}
\begin{tabular}{cl}
  \hline
  Символ & Смысл \\
  \hline
  $H$                    & высота кровли пласта (постоянна), м \\
  $h_g(x,t) = H - h(x,t)$ & толщина газового слоя, м \\
  $k_g$                  & проницаемость для газа, м$^2$ \\
  $\mu_g$                & вязкость газа, Па$\cdot$с \\
  $V_g(t)$               & объём газовой шапки (на единицу длины
                           скважины), м$^2$ \\
  $P_g(x,t)$             & давление в газовой шапке \\
  $\bar{P}_g(t)$         & средневзвешенное давление в газовой шапке \\
  $v_g(x,t)$             & горизонтальная скорость фильтрации газа, м/с \\
  $c_g$                  & сжимаемость газа, Па$^{-1}$ \\
  $\dot{m}_{\text{ГНК}}$ & массовый поток газа через ГНК (дегазация нефти
                           или конденсация), кг/(м$^2\cdot$с) \\
  \hline
\end{tabular}
\end{center}

%-----------------------------------------------------------------------------
\subsection{Подход~A. Материальный баланс газовой шапки (0D)}
\label{sec:gas_cap_0D}
%-----------------------------------------------------------------------------

\paragraph{Физическое обоснование.}
Вязкость газа на 1--2~порядка ниже вязкости нефти
($\mu_g/\mu_o \sim 0{,}01$). Характерное время выравнивания давления
в газовой шапке:
\begin{equation}\label{eq:tau_gas}
  \tau_g \sim \frac{\varphi\,\mu_g\,c_g\,W^2}{k_g}
\end{equation}
на порядки меньше характерного времени движения ГНК
$\tau_o \sim \varphi\,\mu_o\,W^2/(k\,\Delta\rho\,g\,h_0)$.
Поэтому давление в газовой шапке можно считать
\textbf{пространственно однородным}: $P_g(x,t) \approx \bar{P}_g(t)$.

\paragraph{Материальный баланс.}
Объём газовой шапки (на единицу длины скважины в направлении~$y$):
\begin{equation}\label{eq:Vg}
  V_g(t) = \varphi \int_0^{W} \bigl[H - h(x,t)\bigr]\,dx.
\end{equation}

Газ в газовой шапке подчиняется закону реального газа:
$P_g\,V_g = n_g\,Z\,R\,T_g$, где $n_g$~--- количество молей газа.
Дифференцируя:
\begin{equation}\label{eq:gas_cap_balance}
  \boxed{\;
    \frac{d\bar{P}_g}{dt}
    = -\,\frac{\bar{P}_g}{V_g}\,\frac{dV_g}{dt}
    + \frac{\bar{P}_g}{n_g}\,\frac{dn_g}{dt}.
  \;}
\end{equation}

Первое слагаемое~--- изменение давления из-за изменения объёма
(движение ГНК):
\begin{equation}\label{eq:dVg_dt}
  \frac{dV_g}{dt}
    = -\,\varphi \int_0^{W} \frac{\partial h}{\partial t}\,dx.
\end{equation}
Когда $h$ уменьшается (ГНК поднимается, газ расширяется),
$dV_g/dt > 0$, давление падает.

Второе слагаемое~--- поступление газа, выделившегося при дегазации
нефти:
\begin{equation}\label{eq:dng_dt}
  \frac{dn_g}{dt}
    = \frac{\varphi}{M_g} \int_0^{W}
      \frac{\partial}{\partial t}
      \bigl(h\,S_g\,\rho_g\bigr)\,dx,
\end{equation}
где учтён только газ, перешедший из нефтяной фазы в газовую шапку
(при $S_g$ выше критической газонасыщенности газ становится подвижным
и мигрирует вверх).

\paragraph{Поле скорости газа.}
При однородном давлении горизонтальный градиент давления в газе
пренебрежимо мал, и поток газа~--- преимущественно вертикальный.
Скорость вертикальной миграции газа (из нефтяного слоя в газовую шапку)
определяется из баланса массы газового компонента:
\begin{equation}\label{eq:vg_vertical}
  v_{g,z}(x,t) = -\,\frac{k\,k_{r,g}\,\Delta\rho\,g}{\mu_g},
\end{equation}
что представляет собой закон Дарси для вертикального всплытия газа
в поле силы тяжести (при пренебрежении горизонтальным градиентом
давления).

\paragraph{Достоинства и ограничения.}
Подход~A даёт замкнутую модель из одного ODE~\eqref{eq:gas_cap_balance}
для $\bar{P}_g(t)$, не требующую дополнительной сетки в газовой шапке.
Он адекватен при высокой подвижности газа и умеренных размерах
газовой шапки. Ограничение: не описывает пространственную неоднородность
давления в газовой шапке и горизонтальные потоки газа.

%-----------------------------------------------------------------------------
\subsection{Подход~B. Приближение Дюпюи для газовой шапки (PDE)}
\label{sec:gas_cap_PDE}
%-----------------------------------------------------------------------------

\paragraph{Идея.}
Применяем к газовой шапке тот же приём, что Konieczek применил к
нефтяному слою: вертикально~--- гидростатическое равновесие,
горизонтально~--- фильтрация.

\paragraph{Геометрия.}
Газовая шапка ограничена сверху кровлей пласта ($z = H = \mathrm{const}$),
снизу~--- ГНК ($z = h(x,t)$). Толщина газового слоя
$h_g(x,t) = H - h(x,t)$.

\paragraph{Давление в газовой шапке.}
Гидростатика в газе (отсчёт от кровли):
\begin{equation}\label{eq:Pg_hydrostatic}
  P_g(x,z,t) = P_{\text{кров}}(x,t) + \rho_g\,g\,(H - z),
\end{equation}
где $P_{\text{кров}}(x,t)$~--- давление на кровле.
Горизонтальный градиент давления:
\begin{equation}\label{eq:dPg_dx}
  \frac{\partial P_g}{\partial x}
    = \frac{\partial P_{\text{кров}}}{\partial x}
\end{equation}
(не зависит от~$z$, как и в нефти~--- следствие приближения Дюпюи).

\paragraph{Скорость фильтрации газа.}
Горизонтальная скорость:
\begin{equation}\label{eq:vg_horizontal}
  v_g(x,t) = -\,\frac{k_g}{\mu_g}\,
              \frac{\partial P_{\text{кров}}}{\partial x}.
\end{equation}

Горизонтальный расход газа через вертикальное сечение (на единицу
длины скважины):
\begin{equation}\label{eq:qg_horizontal}
  q_g(x,t) = h_g\,v_g
    = -\,\frac{k_g}{\mu_g}\,h_g\,
      \frac{\partial P_{\text{кров}}}{\partial x}.
\end{equation}

\paragraph{Уравнение неразрывности газа.}
Баланс массы газа в слое толщиной~$h_g$ с учётом сжимаемости
и притока газа через ГНК:
\begin{equation}\label{eq:gas_continuity}
  \varphi\,\frac{\partial}{\partial t}(h_g\,\rho_g)
  + \frac{\partial}{\partial x}(h_g\,\rho_g\,v_g) = \Sigma_g,
\end{equation}
где $\Sigma_g$~--- источник массы газа (поступление дегазированного газа
через ГНК из нефтяного слоя, кг/(м$^2\cdot$с)).

Так как $h_g = H - h$, имеем $\partial h_g/\partial t = -\partial h/\partial t$
(из уравнения для нефти~\eqref{eq:gorm_modified}).

\paragraph{Связь плотности и давления.}
Для реального газа $\rho_g = P\,M_g/(Z\,R\,T)$. При умеренных
изменениях давления линеаризуем:
\begin{equation}\label{eq:rho_g_linear}
  \rho_g \approx \rho_{g,0}\,
    \bigl[1 + c_g\,(P_{\text{кров}} - P_0)\bigr],
  \qquad c_g = \frac{1}{P_0}
\end{equation}
(для идеального газа; для реального~--- $c_g = 1/P - Z'/Z$).

\paragraph{Линеаризованное уравнение.}
Подставляя~\eqref{eq:vg_horizontal} и~\eqref{eq:rho_g_linear}
в~\eqref{eq:gas_continuity} и пренебрегая членами второго порядка
малости:
\begin{equation}\label{eq:gas_PDE}
  \boxed{\;
    \varphi\,c_g\,h_g\,\frac{\partial P_{\text{кров}}}{\partial t}
    = \frac{k_g}{\mu_g}\,\frac{\partial}{\partial x}
      \left(h_g\,\frac{\partial P_{\text{кров}}}{\partial x}\right)
    + \frac{\Sigma_g}{\rho_{g,0}}
    + \varphi\,\frac{\partial h}{\partial t}\,
      \frac{P_{\text{кров}}}{P_0}.
  \;}
\end{equation}

Это линейное параболическое PDE для $P_{\text{кров}}(x,t)$
с коэффициентами, зависящими от~$h(x,t)$ (известного из
уравнения~\eqref{eq:gorm_modified}).

Слагаемые в правой части:
\begin{itemize}[nosep]
  \item первое~--- горизонтальная фильтрация газа (диффузия давления);
  \item второе~--- источник газа от дегазации нефти;
  \item третье~--- изменение объёма газовой шапки из-за движения ГНК.
\end{itemize}

\paragraph{Источник газа через ГНК.}
\begin{equation}\label{eq:Sigma_g}
  \Sigma_g(x,t) = \varphi\,\rho_{g,\text{ст}}\,
    \frac{\partial}{\partial t}
    \bigl[h\,S_g\,B_g^{-1}\bigr]_{\text{нефт.\ слой}},
\end{equation}
где производная берётся только по газу, выделившемуся из нефти и
мигрировавшему в газовую шапку.

\paragraph{Граничные условия.}
На скважине ($x = 0$):
\begin{equation}\label{eq:gas_BC_well}
  P_{\text{кров}}(0, t) = P_{\text{заб}}(t)
    + \rho_g\,g\,(H - z_w)
  \quad\text{(если газ прорвался к скважине)},
\end{equation}
или $q_g(0,t) = 0$ до прорыва газа.

На внешнем контуре ($x = W$):
\begin{equation}\label{eq:gas_BC_outer}
  \frac{\partial P_{\text{кров}}}{\partial x}\bigg|_{x=W} = 0
  \quad\text{(непротекание)}.
\end{equation}

\paragraph{Начальное условие.}
\begin{equation}\label{eq:gas_IC}
  P_{\text{кров}}(x, 0) = P_{\text{газ},0}
  \quad\text{(начальное давление газовой шапки)}.
\end{equation}

\paragraph{Поле скорости газа.}
После решения~\eqref{eq:gas_PDE} поле горизонтальной скорости
восстанавливается по~\eqref{eq:vg_horizontal}. Вертикальная
составляющая определяется из уравнения неразрывности:
\begin{equation}\label{eq:vg_z_from_cont}
  v_{g,z}(x,z,t) = -\,\frac{1}{\rho_g}
    \int_z^{H} \frac{\partial(\rho_g\,v_g)}{\partial x}\,dz'
    \approx -\,(H-z)\,\frac{\partial v_g}{\partial x}.
\end{equation}

%-----------------------------------------------------------------------------
\subsection{Сравнение подходов A и B}
%-----------------------------------------------------------------------------

\begin{center}
\renewcommand{\arraystretch}{1.3}
\begin{tabular}{|l|c|c|}
  \hline
  & \textbf{Подход~A (0D)} & \textbf{Подход~B (PDE)} \\
  \hline
  Тип уравнения      & 1 ODE для $\bar{P}_g(t)$
                      & 1 PDE для $P_{\text{кров}}(x,t)$ \\
  $P_g(x)$           & однородное
                      & распределённое \\
  Гориз.\ потоки газа & нет  & да \\
  Верт.\ потоки газа  & по Дарси & из неразрывности \\
  Доп.\ сетка         & нет  & та же, что для $h$ \\
  Сложность/шаг       & $O(1)$ & $O(N)$, прогонка \\
  Когда адекватен     & высокая подвижность газа,
                        $\tau_g \ll \tau_o$
                      & умеренная подвижность,
                        большая газ.\ шапка \\
  \hline
\end{tabular}
\end{center}

%-----------------------------------------------------------------------------
\subsection{Интеграция в модель~6}
%-----------------------------------------------------------------------------

При использовании подхода~A модель~6 остаётся системой из двух PDE
($h$,~$T$) плюс одно ODE ($\bar{P}_g$). Давление в нефтяном слое
определяется по~\eqref{eq:gorm_Pavg} с $P_{\text{газ}} = \bar{P}_g(t)$.

При использовании подхода~B добавляется третье PDE
для $P_{\text{кров}}(x,t)$, решаемое на той же сетке, что
и уравнения для~$h$ и~$T$. Порядок на каждом временном шаге:
\begin{enumerate}[nosep]
  \item Из $h^n$ и $P_{\text{кров}}^n$ вычислить $\bar{P}^n$,
        $R_s^n$, $S_g^n$, $\alpha_{\mathrm{eff}}^n$, $\Sigma_g^n$.
  \item Решить~\eqref{eq:gorm_modified} $\to h^{n+1}$
        (прогонка).
  \item Решить~\eqref{eq:gas_PDE} $\to P_{\text{кров}}^{n+1}$
        (прогонка).
  \item Вычислить $v_g^{n+1}$ по~\eqref{eq:vg_horizontal},
        источники $Q_{\text{дег}}$, $Q_{\mathrm{JT}}$.
  \item Решить~\eqref{eq:gorm_energy} $\to T^{n+1}$
        (прогонка).
\end{enumerate}

Три прогонки за шаг, $O(N)$ операций~--- задача остаётся
вычислительно лёгкой.

%=============================================================================
\section{Сравнение моделей}
%=============================================================================

\begin{center}
\renewcommand{\arraystretch}{1.3}
\begin{tabular}{|l|c|c|c|c|c|c|}
  \hline
  & \textbf{М.\,1} & \textbf{М.\,2} & \textbf{М.\,3}
  & \textbf{М.\,4} & \textbf{М.\,5} & \textbf{М.\,6} \\
  \hline
  Дросселирование    & да  & нет & да  & да  & да  & да \\
  Разгазирование     & нет & да  & да  & да  & да  & да \\
  Нестационарность   & нет & нет & нет & да  & да  & да \\
  Простр.\ распред.  & нет & нет & нет & нет & да  & да \\
  Гидродинамика      & нет & нет & нет & задана & задана & GORM \\
  Газ.\ шапка        & нет & нет & нет & нет & нет & A/B \\
  Тип уравнений      & алг. & алг. & алг. & ОДУ & 1 PDE & 2--3 PDE \\
  Неизвестные        & $\Delta T$ & $\Delta T$ & $T_2$
    & $P,T,S_o$ & $T(r,t)$ & $h, T, [P_g]$ \\
  \hline
\end{tabular}
\end{center}

\paragraph{Рекомендация.}
Для практических оценок оптимальна комбинация Моделей~1 и~2:
\begin{equation}\label{eq:combined}
  \Delta T_{\mathrm{total}} \approx
    \underbrace{\mu_{\mathrm{JT}}\,\Delta P_{\text{выше } P_{\text{нас}}}}_{\text{дросселирование}}
  + \underbrace{\Delta T_{\mathrm{flash}}(P_{\text{нас}} \to P_2)}_{\text{разгазирование}}.
\end{equation}
Это даёт физически обоснованный результат при минимуме входных данных
и может быть реализовано в 10--15~строк кода.
Для детального моделирования с учётом пространственного распределения
и самосогласованной гидродинамики рекомендуется модель~6 (GORM с
дегазацией и температурой).

\begin{thebibliography}{9}

\bibitem{konieczek1990}
  Konieczek,~J.
  \textit{The Concept of Critical Rate in Gas Coning and Its Use in
  Production Forecasting}.
  SPE~20722, presented at the 65th Annual Technical Conference,
  New Orleans, September 1990.

\bibitem{mjaavatten2006}
  Mjaavatten,~A., Aasheim,~R., Saelid,~S., Gronning,~O.
  \textit{A Model for Gas Coning and Rate-Dependent Gas/Oil Ratio in
  an Oil-Rim Reservoir}.
  SPE Reservoir Evaluation \& Engineering, October 2008.
  SPE~102390.

\end{thebibliography}

\end{document}
